% Compile with pdflatex (run twice for the table of contents / \Cref).
\documentclass[11pt]{article}

% evan.sty by Evan Chen (Boost Software License 1.0)
% https://github.com/vEnhance/dotfiles/blob/main/texmf/tex/latex/evan/evan.sty
% Options: sexy = colored boxes for theorems + colored section heads,
%          hints, british, nothm, nofancy, ... (see the top of evan.sty)
\usepackage[sexy]{evan}

\newcommand{\MM}{\mathbb M}

\title{Problem set 2 solution}
\author{Anakawee Chujinda}

% evan.sty hardcodes its author into the running header; reset it.
\lhead{Anakawee Chujinda}
\rhead{\nouppercase{\leftmark}}

\begin{document}
\maketitle
\newpage

\section{Problems Set 2}
\begin{problem}[10pt]
  Let $\FF$ be an ordered field with $1 \neq 0$. Show that $1 > 0$.
  \emph{Hint: show $(-1)(-1) = 1$ first.}
\end{problem}

\begin{problem}[10pt]
  Recall that two rational numbers $\frac{n_1}{m_1}$ and $\frac{n_2}{m_2}$ are the same,
  denoted by $\frac{n_1}{m_1} = \frac{n_2}{m_2}$, if $n_1 m_2 = n_2 m_1$.
  Show that the addition
  \[
    \frac{n}{m} + \frac{p}{q} := \frac{nq + mp}{mq}
  \]
  is well-defined. That is, if $\frac{n_1}{m_1} = \frac{n_2}{m_2}$ and
  $\frac{p_1}{q_1} = \frac{p_2}{q_2}$, then
  \[
    \frac{n_1}{m_1} + \frac{p_1}{q_1} = \frac{n_2}{m_2} + \frac{p_2}{q_2}.
  \]
\end{problem}
\begin{problem}[10pt]
  Find $\sup E$ and $\inf E$ for the following sets $E$:
  \begin{enumerate}
    \item $E = \{ n \in \ZZ \mid n < \sqrt{12} \}$
    \item $E = \{ r \in \QQ \mid r < \sqrt{12} \}$
    \item $E = \{ x \in \RR \mid x^2 - x - 1 < 0 \}$
    \item $E = \left\{ \frac{n^2+n}{n^2+1} \;\middle|\; n \in \NN \right\}$
  \end{enumerate}
  No proof is needed for the answer.
\end{problem}

\begin{problem}[10pt]
  Let $\MM$ be the set of polynomials with integer coefficients,
  \[
    \MM := \{ f(x) = a_0 + a_1 x + a_2 x^2 + \dots + a_n x^n \mid a_i \in \ZZ \}.
  \]
  Define the relation $0 \prec f$ if $0 < f(x)$ for $x$ large enough. To be more precise,
  we say $0 \prec f$ if there exists $M > 0$ such that $f(x) > 0$ for all $x > M$.
  Then define $f \prec g$ if $0 \prec g - f$.
  Show that $(\MM, \prec)$ is an ordered set.

  You can use the following fact directly (without proving it): if
  $f(x) = a_0 + a_1 x + a_2 x^2 + \dots + a_n x^n$ with $0 < a_n$, then $0 < f(x)$ for
  $x$ large enough.
\end{problem}

\begin{problem}[10pt]
  Let $(\MM, \prec)$ be the ordered set defined in Problem 4.
  Show that $(\MM, \prec)$ \textbf{doesn't} satisfy the Archimedean property.
\end{problem}

\newpage
\begin{problem}[10pt]
  The greatest lower bound of a set $E$ is defined to be the number $\beta$ which has the
  following properties:
  \begin{itemize}
    \item For all $x \in E$, $\beta \le x$.
    \item Suppose $\alpha \le x$ for all $x \in E$. Then $\beta \ge \alpha$.
  \end{itemize}
  Show that for any non-empty set $E \subset \RR$ which is bounded from below, $E$ has the
  greatest lower bound.
\end{problem}
\begin{proof}
  Let $F = \{-x \mid x \in E \} \subset \RR$, then $F$ is bounded from above, let $a = \sup F$ $\therefore -x \geq a$ so that $x \geq -a$ for all $x \in E$, then $-a$ is the lower bound of $E$.

  Assume that there are $b$ such that $b$ is the lower bound of $E$. $\therefore x \geq b \Longrightarrow -x \leq -b$ then $-b$ is the upper bound of $F$. Since $a = \sup F$ then $a \leq -b$ then $-a \geq b$. Then $-a$ is the greater upper bound of $E$. therefore, $E$ has largest lower bound.
\end{proof}

\begin{problem}[20pt]
  Show that for all real numbers $x \in \RR$, there exists a real number $y \in \RR$ such
  that $y^3 = x$. \emph{Hint: start with the case $x > 0$.}
\end{problem}
\begin{proof}
  WLOG, we let $x > 0$ Let $A = \{a \subset \RR \mid 0 < a^3 < x\} \subseteq \RR$ since $A$ is bounded, there the there's a $\sup$.
  Let $y = \sup A$, we'll show that $y^3 = x$.

  Assume that $y^3 > x$, choose $h > 0$ such that $0 < h < \dfrac{y^3-x}{3y^2}$. therefore
  \[
  (y-h)^3 = y^3 - 3y^2h + 3yh^2 - h^3 > x
  \]
  so that $y-h$ is upper bound of $A$ but it contradicts that $y$ is smallest.
  
  Assume that $y^3 < x$, we choose $0 < h < 1$ such that $h < \dfrac{x-y^3}{3y^2 + 3y + 1}$ so that $(y+h)^3 < x$ then $y + h \in A$. Therefore, $y > y + h$ which impossible.

  Hence $y^3 = x$
\end{proof}

\end{document}
